% SC2005 Ch06 — Memory Management (no \documentclass)
\chapter{Memory Management}
\label{ch:memory}

\begin{quote}
\emph{Memory is the workspace of computation.} How the OS partitions, translates,
and protects it determines both performance and isolation.
\end{quote}

\noindent\textbf{Learning Outcomes.} After this chapter you will be able to:
\begin{enumerate}[label=(L\arabic*)]
  \item explain address binding, logical vs.\ physical addresses;
  \item evaluate contiguous allocation, fragmentation, and compaction;
  \item describe paging, TLBs, and multi-level page tables;
  \item compare paging vs.\ segmentation and segmented paging;
  \item reason about protection, sharing, and page-table overhead.
\end{enumerate}

% ============================================================
\section{Address Binding and Spaces}
\label{sec:binding}

A program's addresses are \emph{logical} (virtual); hardware executes on
\emph{physical} addresses after translation via the MMU. Binding can occur at
compile time, load time, or execution time (the most flexible, requiring base +
limit or paging hardware).

Logical addresses from the CPU go through relocation (add base) or page-table
lookup before reaching RAM. The user never sees physical addresses.

\subsection{Contiguous Allocation}
Early systems placed each process in one contiguous block $[base, base+limit)$.
Hardware checks every access: $base \le addr < base+limit$. Simple and fast,
but causes \emph{external fragmentation} --- free memory is scattered in holes
too small for new processes. \emph{Compaction} (moving processes) is expensive.

Allocation strategies for holes:
\begin{itemize}
  \item \textbf{First-fit:} use first hole large enough --- fast, low overhead.
  \item \textbf{Best-fit:} smallest sufficient hole --- reduces waste but leaves tiny slivers.
  \item \textbf{Worst-fit:} largest hole --- tends to leave usable fragments.
\end{itemize}

\begin{figure}[h!]
\centering
\begin{tikzpicture}[scale=0.74, every node/.style={font=\scriptsize},
  m/.style={draw, minimum width=2.4cm, align=center}]
  \draw[thick] (0,0) rectangle (2.4,3.2);
  \node[m, fill=blue!14, minimum height=0.7cm] at (1.2,2.75) {OS};
  \node[m, fill=green!14, minimum height=0.8cm] at (1.2,2.05) {$P_1$};
  \node[m, fill=orange!14, minimum height=0.4cm] at (1.2,1.45) {Hole 200\,KB};
  \node[m, fill=violet!12, minimum height=1.0cm] at (1.2,0.75) {$P_2$};
  \node[m, fill=gray!14, minimum height=0.25cm] at (1.2,0.12) {Hole 80\,KB};
  \node[font=\tiny, align=left] at (4.2,1.5) {External\\fragmentation:\\free scattered\\in small holes};
  \draw[-{Stealth}, thick, red!60!black] (2.5,1.45) -- (3.3,1.45);
  \node[font=\tiny, fill=white] at (1.2,-0.35) {Physical memory (contiguous)};
\end{tikzpicture}
\caption{Contiguous allocation with external fragmentation.}
\label{fig:contiguous}
\end{figure}

% ============================================================
\section{Paging}
\label{sec:paging}

\emph{Paging} eliminates external fragmentation by dividing logical memory into
fixed-size \emph{pages} (e.g., 4\,KB) and physical memory into \emph{frames}
of the same size. A \emph{page table} maps page number $p$ to frame $f$; offset
$d$ is preserved: physical address $= f \times \text{PAGE\_SIZE} + d$.

\begin{figure}[h!]
\centering
\begin{tikzpicture}[scale=0.74, every node/.style={font=\scriptsize},
  b/.style={draw, rounded corners, align=center, minimum width=2.2cm, minimum height=0.55cm}]
  \node[b, fill=blue!12] (vpn) at (0,1.8) {Virtual addr\\ $p$ \,|\, $d$};
  \node[b, fill=orange!16] (pt) at (3.2,1.8) {Page table\\ $p \to f$};
  \node[b, fill=green!14] (phys) at (1.6,0.4) {Physical addr\\ $f$ \,|\, $d$};
  \node[b, fill=violet!12] (mem) at (1.6,-0.75) {RAM frame $f$};
  \draw[-{Stealth}, thick, blue!60!black] (vpn) -- (pt) node[midway,above,font=\tiny]{lookup};
  \draw[-{Stealth}, thick, blue!60!black] (pt) -- (phys);
  \draw[-{Stealth}, thick] (phys) -- (mem);
  \draw[-{Stealth}, thick, dashed, red!60!black] (vpn) -- (phys) node[midway,left,font=\tiny]{$d$ preserved};
\end{tikzpicture}
\caption{Paging address translation: page number maps to frame, offset passes through.}
\label{fig:paging}
\end{figure}

\subsection{TLB}
Every memory access would need a page-table lookup without caching. The
\emph{Translation Lookaside Buffer} (TLB) is a small associative cache of
recent $p\to f$ mappings. On a TLB miss, hardware or OS walks the page table
and refills the TLB. Effective access time:
\[
  t_{\text{eff}} = h \cdot (t_{\text{TLB}}+t_{\text{mem}}) + (1-h)\cdot(t_{\text{TLB}}+2t_{\text{mem}})
\]
with hit rate $h$.

\begin{example}[TLB effect]
With $h=0.98$, $t_{\text{TLB}}=5$\,ns, $t_{\text{mem}}=80$\,ns:
$t_{\text{eff}}=0.98\cdot85+0.02\cdot165=86.6$\,ns vs.\ $160$\,ns without TLB.
Even a 2\% miss rate doubles cost without the TLB.
\end{example}

\subsection{Multi-Level Page Tables}
A flat page table for 32-bit addresses with 4\,KB pages has $2^{20}$ entries
($\sim$4\,MB per process). \emph{Two-level} paging splits $p$ into
$p_1\,|\,p_2$: the outer table points to inner tables allocated only when
needed, drastically reducing overhead for sparse address spaces. 64-bit systems
use 3--4 levels.

% ============================================================
\section{Segmentation and Combined Schemes}
\label{sec:segmentation}

\emph{Segmentation} divides memory logically (code, data, stack, heap) into
variable-size segments with base/limit per segment. It matches programmer
view and enables sharing/protection per segment, but reintroduces external
fragmentation.

\textbf{Segmented paging:} each segment has its own page table --- combining
logical organization with paging's freedom from external fragmentation. Intel
x86 historically used this hybrid.

\begin{figure}[h!]
\centering
\begin{tikzpicture}[scale=0.74, every node/.style={font=\scriptsize},
  s/.style={draw, rounded corners, align=center, minimum width=2.0cm, minimum height=0.6cm}]
  \node[s, fill=blue!12] (va) at (0,1.6) {Virtual addr\\ $s$ \,|\, $p$ \,|\, $d$};
  \node[s, fill=orange!15] (seg) at (3.0,1.6) {Segment table\\ $s\to$ page table};
  \node[s, fill=green!14] (pt) at (3.0,0.4) {Page table\\ $p\to f$};
  \node[s, fill=violet!12] (pa) at (0,0.4) {Physical addr\\ $f$ \,|\, $d$};
  \draw[-{Stealth}, thick, blue!60!black] (va)--(seg);
  \draw[-{Stealth}, thick, blue!60!black] (seg)--(pt);
  \draw[-{Stealth}, thick, blue!60!black] (pt)--(pa);
\end{tikzpicture}
\caption{Segmented paging: segment $\to$ page table $\to$ frame.}
\label{fig:segpage}
\end{figure}

Paging suffers \emph{internal fragmentation} (last page partially unused,
average $\frac{1}{2}$ page per segment). Protection bits (R/W/X, user/supervisor)
and sharing (shared frames, copy-on-write) are enforced per page/segment.

% ============================================================
\begin{remark}[Internal vs.\ external fragmentation]
Paging has no external fragmentation (any free frame fits any page) but pays
internal fragmentation: the last page of each process wastes on average
$\text{PAGE\_SIZE}/2$ bytes. With 4\,KB pages and many small processes this
adds up, motivating variable page sizes (huge pages: 2\,MB/1\,GB) for
large mappings.
\end{remark}

\begin{lstlisting}[language=C,caption={Page-table walk (two-level, sketch)}]
// va = p1 | p2 | offset  (10 | 10 | 12 for 4KB pages, 32-bit)
pte_t *pt = outer[p1];          // outer table lookup
if (!pt) fault("outer miss");
frame = pt[p2].frame;           // inner table lookup
phys = (frame << 12) | offset;
\end{lstlisting}

% TLB shootdown is a key multicore overhead
% covered further in virtual memory chapter
\begin{remark}[Why TLB shootdown matters]
On multicore, updating a page-table entry requires invalidating stale TLB
entries on other cores (TLB shootdown via IPIs). Frequent \texttt{mmap/mprotect}
can thus be surprisingly expensive on large machines.
\end{remark}

\begin{example}[Two-level page-table size]
For a 32-bit VA with 4\,KB pages, $p_1$ and $p_2$ are 10 bits each. The outer
table has $1024$ entries; each inner table has $1024$ entries. A process using
only two far-apart pages needs just 1 outer entry $+$ 2 inner tables ($\sim$12\,KB)
instead of a flat 4\,MB table --- a $300\times$ saving.
\end{example}

\begin{remark}[Huge pages]
Using 2\,MB or 1\,GB huge pages reduces TLB pressure (one entry covers more)
at the cost of higher internal fragmentation. Databases and HPC workloads
often enable huge pages for this reason.
\end{remark}

\section{Exercises}
\label{sec:ch06-ex}
\begin{exercise}
Given logical address 0x1234 with page size 1\,KB (10-bit offset), find page
number and offset. If page 4 maps to frame 7, what is the physical address?
\end{exercise}
\begin{exercise}
For $t_{\text{mem}}=100$\,ns, $t_{\text{TLB}}=10$\,ns, and hit rate 95\%,
compute effective access time with and without TLB.
\end{exercise}
\begin{exercise}
Compare best-fit vs.\ first-fit allocation on a sequence of allocations and
frees. Which fragments more and why?
\end{exercise}
\begin{exercise}
Why do 64-bit systems need multi-level page tables? Estimate table size for a
flat table with 4\,KB pages.
\end{exercise}
\begin{exercise}
Explain how copy-on-write for \texttt{fork} uses page-table protection bits to
defer copying until a write occurs.
\end{exercise}
